\uuid{PgOg}
\titre{Racine}
\auteur{bodin ; hanani ; mzari}
\organisation{exo7}
\licence{CC BY-NC-SA 4.0}
\datecreate{2020-01-20}
\niveau{L1}
\module{Algèbre}
\chapitre{Polynôme, fraction rationnelle}
\sousChapitre{Racine, décomposition en facteurs irréductibles}
\difficulte{2}
\theme{Racines, factorisation, racine}
\exo7qcmid{lille-1.polynomes.022}

\contenu{
\begin{question}
\qclassification{105.03}
\qtags{motcle=racine}
Soit $P(X) = X^6 + 4 X^5 + X^4 - 10 X^3 - 4 X^2 + 8 X$.
\begin{answers}
    \bad{$-1$ est une racine double.}

    \bad{$0$ est une racine double.}

    \good{$1$ est une racine double.}

    \bad{$-2$ est une racine double.}
\end{answers}
\begin{explanations}
Pour une racine double il faut $P(a)=0$, $P'(a)=0$ et $P''(a)\neq0$.
En fait $P(X) = X(X+2)^3(X-1)^2$.
\end{explanations}
\end{question}
}
